\section[Cote 156-2 : les rives d'une partie fermée]{Cote 156-2 : les rives d'une partie fermée\protect\verifie{Folder156_2}}\label{sec:156-2}

La cote 156-2, chapitre~II de \emph{Vers une géométrie des formes}, dix-huit pages des 6 et 7~juin 1986 intitulées \q{Réalisations topologiques des réseaux}, cherche à couper un espace quelconque le long d'une partie fermée sans le supposer variété ni triangulé. Pour $A$ fermé dans $X$, $j : X \smallsetminus A \hookrightarrow X$, et $\mathrm{Comp}(Z)$ l'anneau booléen des parties ouvertes et fermées de $Z$, la page~14 pose $\mathrm{Riv}(A, X) = \varinjlim_{V \supset A} \mathrm{Comp}(V \smallsetminus A)$, limite sur les voisinages de $A$, et $\widetilde A = \mathrm{Spec}(j_*\mathbf F_2)$ ; la page~15 en tire une application canonique $(*)$ $\mathrm{Riv}(A, X) \to \mathrm{Comp}(\widetilde A)$, \q{injective, et} ajouté au-dessus de \q{bijective si $X$ paracompact} (\lot{156-2}{1}{14}, \lot{156-2}{1}{15}). La lecture prend $\widetilde A \to A$ pour le spectre de Stone relatif de $(j_*\mathbf F_2)|_A$ et note que l'injectivité \q{n'utilise rien}.

On code une partie ouverte et fermée de $V \smallsetminus A$ par son indicatrice, fonction $X \to \mathbf F_2$ localement constante sur $V \smallsetminus A$ ; $\mathrm{Riv}(A, X)$ est l'anneau de leurs germes le long de $A$, la fibre $B_x$ de $j_*\mathbf F_2$ en $x$ celui de leurs germes en $x$ dans $X \smallsetminus A$. On prend pour $\widetilde A$ l'ensemble des couples $(x, \chi)$, $x \in A$, $\chi : B_x \to \mathbf F_2$ un morphisme d'anneaux (un point du spectre de $B_x$), muni de la topologie engendrée par les parties $\{(y, \chi) : y \in W,\ \chi(f_y) = 1\}$, $W$ ouvert et $f$ localement constante sur $W \smallsetminus A$ ; $(*)$ envoie $\rho$ sur $\{(x, \chi) : \chi(\rho_x) = 1\}$.

\begin{theoreme}[D'après les pages 14 et 15, injectivité\piste{156-2-rives-cut-along-closed-set}]
Pour toute partie $A$ de tout espace $X$, l'image de chaque rive par $(*)$ est ouverte et fermée dans $\widetilde A$, $(*)$ est un morphisme d'algèbres de Boole, et $(*)$ est injective. L'anneau $\mathrm{Riv}(A, X)$ est nul si et seulement si $A$ est ouvert.
\end{theoreme}

\noindent L'énoncé tient, sans hypothèse : voir les théorèmes \texttt{toComp\_injective}, \texttt{toComp\_mul}, \texttt{toComp\_one}, \texttt{toSpecSet\_compl} et \texttt{riv\_one\_eq\_zero\_iff} dans \href{https://github.com/Commutator-IO/grothendieck-project/blob/main/lean/Grothendieck/Folder156_2.lean}{\texttt{Folder156\_2.lean}}. La fermeture de $A$ ne sert pas. Deux germes distincts en $x$ sont séparés par un point du spectre de $B_x$, obtenu ici comme évaluation le long d'un ultrafiltre plus fin que le filtre des voisinages de $x$ dans $X \smallsetminus A$. La dernière assertion est la correction \q{$A$ non ouvert} que la lecture apporte au \q{$A \neq \emptyset$} de la remarque sur les bi-rives. Ne sont pas démontrés : la bijectivité pour $X$ paracompact, l'identification $\mathrm{Comp}(\widetilde A) \simeq \Gamma(A, (j_*\mathbf F_2)|_A)$, le cas régulièrement immergé et l'espace découpé $\overline X$ des pages~16 et~17.
